\subsection{Теорема об инвариантности ранга при элементарных преобразованиях. Нахождение ранга матрицы методом элементарных преобразований}

\begin{definition}
Рангом матрицы $A$ называется размер максимального отличного от нуля минора этой матрицы. Второе определение: ранг матрицы — это максимальное количество её линейно независимых строк (или столбцов).
\end{definition}

\begin{theorem}[Об инвариантности ранга]
При элементарных преобразованиях ранг матрицы не меняется.
\end{theorem}

\begin{tcolorbox}[title=Доказательство, breakable]
Рассмотрим три типа элементарных преобразований для строк (для столбцов доказательство аналогично в силу равноправия строк и столбцов):

\textbf{1. Перестановка двух строк.}
Воспользуемся вторым определением ранга. При перестановке строк сами векторы-строки не меняются, меняется лишь их порядок в таблице. Следовательно, максимальное количество линейно независимых строк остаётся прежним. Ранг не меняется.

\textbf{2. Умножение строки на число $\lambda \neq 0$.}
Пусть ранг матрицы $A$ равен $r$. Это значит, что существует минор $M_r \neq 0$, а все миноры $M_{r+1} = 0$. При умножении $i$-й строки на $\lambda$:
\begin{itemize}
    \item Если минор $M_r$ содержит $i$-ю строку, то новый минор $M'_r = \lambda M_r$. Так как $\lambda \neq 0$ и $M_r \neq 0$, то $M'_r \neq 0$.
    \item Если минор $M_r$ не содержит $i$-ю строку, то $M'_r = M_r \neq 0$.
\end{itemize}
В обоих случаях в новой матрице существует минор порядка $r$, не равный нулю. Аналогично, любой минор порядка $r+1$ либо не изменится, либо умножится на $\lambda$, оставаясь равным нулю. По определению, ранг остался равен $r$.

\textbf{3. Прибавление к одной строке другой строки, умноженной на число $\beta$.}
Докажем, что при таком преобразовании ранг не увеличивается. Пусть к первой строке прибавили вторую, умноженную на $\beta$. Рассмотрим минор $M_{r+1}$ новой матрицы. По свойству линейности определителя этот минор можно представить как сумму двух миноров исходной матрицы:
\[ M'_{r+1} = M_{r+1}^{(1)} + \beta M_{r+1}^{(2)} \]
где $M_{r+1}^{(1)}$ — минор исходной матрицы, а $M_{r+1}^{(2)}$ — минор, в котором две строки пропорциональны (или совпадают), если в него вошли обе преобразуемые строки, либо это также минор исходной матрицы. Так как ранг исходной матрицы $r$, все миноры порядка $r+1$ равны нулю. Значит, $M'_{r+1} = 0 + \beta \cdot 0 = 0$. 
Таким образом, ранг новой матрицы $\mathrm{rk}(A') \le \mathrm{rk}(A)$. Поскольку элементарные преобразования обратимы, исходную матрицу можно получить из новой обратным преобразованием того же типа. Следовательно, $\mathrm{rk}(A) \le \mathrm{rk}(A')$. Из двух неравенств получаем $\mathrm{rk}(A') = \mathrm{rk}(A)$.
\end{tcolorbox}

\begin{definition}
Матрица $C$ называется трапециевидной (или ступенчатой), если все её нулевые строки (если они есть) стоят последними, а в ненулевых строках $C_1, C_2, \dots, C_r$ первые ненулевые элементы (лидеры строк) смещаются в каждой следующей строке вправо.
\[
C = \begin{pmatrix}
1 & \dots & \dots & \dots \\
0 & 1 & \dots & \dots \\
\dots & \dots & \dots & \dots \\
0 & 0 & 0 & 1 & \dots \\
0 & 0 & 0 & 0 & 0 \\
\dots & \dots & \dots & \dots & \dots
\end{pmatrix}
\]
\end{definition}

\begin{theorem}
Ранг трапециевидной матрицы равен количеству её ненулевых строк.
\end{theorem}

\begin{tcolorbox}[title=Доказательство, breakable]
Пусть в трапециевидной матрице $r$ ненулевых строк. 
1. Рассмотрим минор порядка $r$, стоящий на пересечении этих $r$ строк и столбцов, в которых расположены их лидеры. После перестановки столбцов (если необходимо) этот минор становится треугольным с ненулевыми элементами на главной диагонали. Его определитель равен произведению диагональных элементов и не равен нулю. Значит, $\mathrm{rk}(C) \ge r$.
2. Любой минор порядка $r+1$ обязательно будет содержать хотя бы одну нулевую строку (так как ненулевых строк всего $r$). Определитель такого минора равен нулю.
Следовательно, ранг равен в точности $r$.
\end{tcolorbox}

\textbf{Нахождение ранга матрицы (алгоритм):}
Для нахождения ранга произвольной ненулевой матрицы $A$:
\begin{enumerate}
    \item С помощью элементарных преобразований строк (перестановка, умножение на число, прибавление строки к строке) и, если удобно, перестановки столбцов, матрица приводится к трапециевидной форме.
    \item В силу доказанной теоремы об инвариантности, ранг при этом не меняется.
    \item Ранг исходной матрицы будет равен числу ненулевых строк полученной трапециевидной матрицы.
\end{enumerate}

\begin{tcolorbox}[title=Источники, colback=gray!10]
\begin{itemize}
  \item Лекция №\,1, тема «Элементарные преобразования и ранг», временной отрезок 04:00--36:00
  \item Лекция №\,2, тема «Метод Гаусса и ранг», временной отрезок 30:00--33:00
  \item PDF «Линейная алгебра\_compressed (1).pdf», раздел 2 «Вычисление ранга матрицы», стр. 3--5
  \item PDF «Вопросы», вопрос №\,2
\end{itemize}
\end{tcolorbox}
